% ============ 第3章 §3.4 关于谱综合----Varopoulos 的方法(原书 pp.249-265) ============
\section{关于谱综合----Varopoulos 的方法}\label{sec:3.4}
\index{谱综合!非谱综合集}
\index{Helson 集}
\index{相对零化}
\index{受限锥性质}
\index{曲率}
\index{第二基本矩阵}
\index{Katznelson 定理}
\index{Wiener--Lévy 定理}
\index{作用函数!(函数空间的)}
\index{作用!(函数空间的情形)}


Schwartz \cite{ch3:49} 证明了 $\mathbf{R}^{3}$ 的单位球面 $S^{2}$ 不是 $L^{1}(\mathbf{R}^{3})$ 中的谱综合集.这里,我们叙述利用这一结果,关于交换 Banach 代数谱综合的 Varopoulos 的研究 \cite{ch3:52}(以及其参考文献 cf. \cite{ch3:23}).此外,也对 L. Schwartz 这一研究在几何方向上的发展作一补充.

\subsection{Schwartz 的结果}\label{sec:3.4.1}

下面叙述的是 Schwartz \cite{ch3:49} 的结果.

\begin{theorem}\label{thm:3.4.1}
当 $n \geq 3$ 时,$S^{n-1}$ 不是 $A(\mathbf{R}^{n})$ 的谱综合集.
\end{theorem}

\begin{proof}
令
\begin{equation*}
I(S^{n-1}) = \{f \in A(\mathbf{R}^{n}): f(y) = 0 \ (y \in S^{n-1})\},
\end{equation*}
\begin{equation*}
I_{1}(S^{n-1}) = \left\{f \in I(S^{n-1}) \cap C^{\infty}(\mathbf{R}^{n}): \frac{\partial f}{\partial y_{1}}(y) = 0 \ (y \in S^{n-1})\right\}.
\end{equation*}
以 $\mu$ 表示 $S^{n-1}$ 上的面积元素,并设 $\mu(S^{n-1}) = 1$.此时
\begin{equation*}
\widehat{\mu}(y) = \int_{S^{n-1}} e^{-iy \cdot x}\dd\mu(x) = \Gamma\left(\frac{n}{2}\right)\left(\frac{2}{|y|}\right)^{\frac{n-2}{2}} J_{\frac{n-2}{2}}(|y|) \quad (n \geq 3,\ y \neq 0)
\end{equation*}
(cf. \cite{ch3:11}).其中,$\Gamma$ 是 Gamma 函数,$J_{\alpha}$ 是第一类 Bessel 函数.此时,对 $y = (y_{1}, \ldots, y_{n})$,
\begin{equation*}
|y_{1}||\widehat{\mu}(y)| \leq C\left(\frac{1}{|y|}\right)^{\frac{n-2}{2}}\left|J_{\frac{n-2}{2}}(|y|)\right||y|
\end{equation*}
由 Bessel 函数的性质($C$ 是适当的常数)
\begin{equation*}
\left\{\begin{array}{ll}
\left|J_{\frac{n-2}{2}}(|y|)\right| \leq C'|y|^{\frac{n-2}{2}} & (|y| \to 0),\\
\left|J_{\frac{n-2}{2}}(|y|)\right| \leq C'|y|^{-\frac{1}{2}} & (|y| \to \infty)
\end{array}\right.
\end{equation*}
可知当 $n \geq 3$ 时,对适当的常数 $C$ 有 $|y_{1}\widehat{\mu}(y)| \leq C$($y \in \mathbf{R}^{n}$).于是,令
\begin{align*}
Tf &= -i\int_{\mathbf{R}^{n}}\tilde{f}(y)y_{1}\widehat{\mu}(y)\frac{\dd y}{2\pi}\\
&\quad \left(\tilde{f} \in C_{c}^{\infty}(\mathbf{R}^{n}) \subset A(\mathbf{R}^{n}),\ f \text{ 是 } \tilde{f} \text{ 的 Fourier 变换}\right)
\end{align*}
则
\begin{equation*}
|Tf| \leq C\int_{\mathbf{R}^{n}}|\tilde{f}(y)|\dd y = \|f\|_{A(\mathbf{R}^{n})}.
\end{equation*}
又因 $C_{c}^{\infty}(\mathbf{R}^{n})$ 在 $A(\mathbf{R}^{n})$ 中稠密,故可得 $|Tf| \leq C\|f\|_{A(\mathbf{R}^{n})}$.进一步,
\begin{align*}
Tf &= -i\int_{\mathbf{R}^{n}}\tilde{f}(y)y_{1}\widehat{\mu}(y)\frac{\dd y}{2\pi}\\
&= \iint (-iy_{1})\tilde{f}(y)e^{-iy \cdot x}\frac{\dd y}{2\pi}\dd\mu(x)\\
&= \int \lim_{h \to 0}\frac{\int \tilde{f}(y)e^{-i(x+h_{1})\cdot y}\frac{\dd y}{2\pi} - \int \tilde{f}(y)e^{-ix \cdot y}\frac{\dd y}{2\pi}}{h}\dd\mu(x)
\end{align*}
\begin{equation*}
= \int \frac{\partial f}{\partial x_{1}}(x)\dd\mu(x)
\end{equation*}
其中设 $h_{1} = (h, 0, \ldots, 0)$.此时,若 $f \in I_{1}(S^{n-1})$,则 $Tf = 0$.可是,在 $1/2 < |x| < 3/2$ 上取满足 $g(x) = x_{1}(|x|^{2} - 1)$ 的 $g \in C_{c}^{\infty}(\mathbf{R}^{n})$,则 $g \in I(S^{n-1})$,而且
\begin{equation*}
Tg = \int_{S^{n-1}} 2x_{1}^{2}\dd\mu(x) > 0.
\end{equation*}
于是由泛函分析的基础知识
\begin{equation*}
I(S^{n-1}) \supsetneqq \overline{I_{1}(S^{n-1})} \supseteq J(S^{n-1})
\end{equation*}
成立.因此,考虑到 $\overline{I_{1}(S^{n-1})}$ 是 $A(\mathbf{R}^{n})$ 的闭理想,$S^{n-1}$ 不是谱综合集.
\end{proof}

\subsection{\texorpdfstring{$A(\mathbf{T}^3)$ 的谱综合}{A(T^3) 的谱综合}}\label{sec:3.4.2}

由 3.4.1 小节可知 $A(\mathbf{R}^{3})$ 是不可谱综合的交换 Banach 代数.利用这一点,可如下证明 $A(\mathbf{T}^{3})$ 不可谱综合.

\begin{proposition}\label{prop:3.4.2}
设
\begin{equation*}
E = \left\{(x_{1}, x_{2}, x_{3}) \in \mathbf{R}^{3}: |x_{1}| \geq \frac{\pi}{2} \text{ 或 } |x_{2}| \geq \frac{\pi}{2} \text{ 或 } |x_{3}| \geq \frac{\pi}{2}\right\},
\end{equation*}
则 $E$ 是 $A(\mathbf{R}^{3})$ 的谱综合集.
\end{proposition}

\begin{proof}
设 $f \in A(\mathbf{R}^{3})$,$f = 0$ on $E$,并令 $f_{\alpha}(x) = \alpha^{3}f(\alpha x)$,$0 < \alpha < 1$,则
\begin{equation*}
\|f_{\alpha} - f\|_{A(\mathbf{R}^{3})} \to 0 \quad (\alpha \to 1),
\end{equation*}
且 $f_{\alpha}$ 在 $E$ 的邻域上为 0.因此 $E$ 是谱综合集.
\end{proof}

\begin{proposition}\label{prop:3.4.3}
在命题~\ref{prop:3.4.2} 的记号下,$E \cup S^{2}$ 不是 $A(\mathbf{R}^{3})$ 的谱综合集.
\end{proposition}

\begin{proof}
假设 $E \cup S^{2}$ 是谱综合集,取 $f \in I(S^{2})$.对 $\varepsilon > 0$,取 $h \in A(\mathbf{R}^{3})$ 使得 $\supp h$ 是紧集且 $\|f - hf\|_{A(\mathbf{R}^{3})} < \varepsilon$.再取 $g \in A(\mathbf{R}^{3})$ 使得 $g$ 在 $S^{2}$ 的邻域上为 0,在 $\supp h \cap E$ 上为 1.此时因 $fh - ghf \in I(E \cup S^{2})$,由 $E \cup S^{2}$ 是谱综合集的假设,存在在 $E \cup S^{2}$ 的邻域上为 0 的 $k \in A(\mathbf{R}^{3})$ 使得 $\|k - (fh - fhg)\|_{A(\mathbf{R}^{3})} < \varepsilon$.从而
\begin{equation*}
\|f - (k + fhg)\|_{A(\mathbf{R}^{3})} \leq \|f - fh\|_{A(\mathbf{R}^{3})} + \|fh - (k + fhg)\|_{A(\mathbf{R}^{3})} < 2\varepsilon,
\end{equation*}
于是 $S^{2}$ 成为 $A(\mathbf{R}^{3})$ 的谱综合集,产生矛盾.
\end{proof}

\begin{proposition}\label{prop:3.4.4}
设 $f \in A(\mathbf{T}^{3})$ 满足 $f(x_{1}, x_{2}, x_{3}) = 0$($\frac{\pi}{2} \leq |x_{1}| \leq \pi$ 或 $\frac{\pi}{2} \leq |x_{2}| \leq \pi$ 或 $\frac{\pi}{2} \leq |x_{3}| \leq \pi$).此时,令
\begin{equation*}
g(x_{1}, x_{2}, x_{3}) = \left\{\begin{array}{ll}
f(x_{1}, x_{2}, x_{3}) & (|x_{j}| \leq \frac{\pi}{2},\ j = 1, 2, 3),\\
0 & (\text{其他})
\end{array}\right.
\end{equation*}
则 $g \in A(\mathbf{R}^{3})$,且存在常数 $C_{1}, C_{2} > 0$ 使得
\begin{equation*}
C_{1}\|g\|_{A(\mathbf{R}^{3})} \leq \|f\|_{A(\mathbf{T}^{3})} \leq C_{2}\|g\|_{A(\mathbf{R}^{3})}.
\end{equation*}
\end{proposition}

\begin{proof}
设 $E' = \{(x_{1}, x_{2}, x_{3}): |x_{j}| \leq \frac{\pi}{2},\ j = 1, 2, 3\}$.取 $h \in C_{c}^{\infty}(\mathbf{R}^{3})$ 使得 $h$ 在 $E'$ 的邻域上为 1 且 $\supp h \subset \{(x_{1}, x_{2}, x_{3}): |x_{j}| \leq \frac{2\pi}{3},\ j = 1, 2, 3\}$.对 $f \in A(\mathbf{T}^{3})$,由 $g(x) = f(x)h(x)$($x = (x_{1}, x_{2}, x_{3})$)得 $g \in A(\mathbf{R}^{3})$ 且
\begin{equation*}
\|g\|_{A(\mathbf{R}^{3})} \leq \|h\|_{A(\mathbf{R}^{3})}\|f\|_{A(\mathbf{T}^{3})}.
\end{equation*}
反之,若 $g \in A(\mathbf{R}^{3})$,则当 $f = gh$ 时,
\begin{equation*}
\widehat{f}(n) = \frac{1}{(2\pi)^{3}}\int_{\mathbf{R}^{3}}g(x)h(x)e^{-ix \cdot n}\dd x \quad (n \in \mathbf{Z}^{3}).
\end{equation*}
由于 $\widehat{h} \in L^{1}(\mathbf{R}^{3})$,
\begin{equation*}
\widehat{f}(n) = \int_{\mathbf{R}^{3}}\widehat{g}(y)\widehat{h}(n - y)\dd y.
\end{equation*}
此时由 $h \in C_{c}^{\infty}(\mathbf{R}^{3})$,存在 $b > 0$ 使得
\begin{equation*}
\sum_{n \in \mathbf{Z}^{3}}|\widehat{h}(n - y)| < b \quad (y \in \mathbf{R}^{3}),
\end{equation*}
故得 $\|f\|_{A(\mathbf{T}^{3})} \leq b\|g\|_{A(\mathbf{R}^{3})}$.
\end{proof}

\begin{theorem}\label{thm:3.4.5}
$A(\mathbf{T}^{3})$ 不可谱综合.
\end{theorem}

\begin{proof}
利用命题~\ref{prop:3.4.2}～3.4.4 的记号.令
\begin{equation*}
\widetilde{E} = \left\{x \in \mathbf{T}^{3}: |x_{1}| \geq \frac{\pi}{2} \text{ 或 } |x_{2}| \geq \frac{\pi}{2} \text{ 或 } |x_{3}| \geq \frac{\pi}{2}\right\}.
\end{equation*}
此时 $\widetilde{E}$ 是 $A(\mathbf{T}^{3})$ 的谱综合集.事实上,若 $h \in I(\widetilde{E}) = \{f \in A(\mathbf{T}^{3}): f = 0 \text{ on } \widetilde{E}\}$,则由命题~\ref{prop:3.4.4},取 $h$ 所对应的 $g \in A(\mathbf{R}^{3})$,则 $g \in I(E)$.由命题~\ref{prop:3.4.2},对 $\varepsilon > 0$,存在在 $E$ 的邻域上为 0 且 $\|g - g_{0}\|_{A(\mathbf{R}^{3})} < \varepsilon$ 的 $g_{0} \in A(\mathbf{R}^{3})$.此时,设 $h_{0}$ 为命题~\ref{prop:3.4.4} 中对应于 $g_{0}$ 的函数,则 $h_{0}$ 在 $\widetilde{E}$ 的邻域上为 0 且满足 $\|h - h_{0}\|_{A(\mathbf{T}^{3})} < C\varepsilon$,所以 $\widetilde{E}$ 是 $A(\mathbf{T}^{3})$ 的谱综合集.另一方面,设 $\widetilde{S^{2}} = \{x = (x_{1}, x_{2}, x_{3}) \in \mathbf{T}^{3}: x_{1}^{2} + x_{2}^{2} + x_{3}^{2} = 1\}$,$\widetilde{K} = \widetilde{E} \cup \widetilde{S^{2}}$,若 $\widetilde{K}$ 成为 $A(\mathbf{T}^{3})$ 的谱综合集,则由命题~\ref{prop:3.4.4},$K = E \cup S^{2}$ 成为 $A(\mathbf{R}^{3})$ 的谱综合集,与命题~\ref{prop:3.4.3} 矛盾.因此 $\widetilde{K}$ 不是 $A(\mathbf{T}^{3})$ 的谱综合集.故 $A(\mathbf{T}^{3})$ 不可谱综合.
\end{proof}

\subsection{\texorpdfstring{交换 Banach 代数 $V(K)$}{交换 Banach 代数 V(K)}}\label{sec:3.4.3}

基于 Varopoulos \cite{ch3:52},设 $K_{1}, K_{2}$ 为紧 Hausdorff 空间,如下定义 $V(K_{1} \times K_{2})$.

\begin{definition}\label{def:3.4.6}
称 $\varphi \in V(K_{1} \times K_{2})$,是指 $\varphi(x, y) = \sum_{j=1}^{\infty} f_{j}(x)g_{j}(y)$,其中 $f_{j} \in C(K_{1})$,$g_{j} \in C(K_{2})$ 满足 $\sum_{j=1}^{\infty}\|f_{j}\|_{\infty}\|g_{j}\|_{\infty} < \infty$.不过,$\|\cdot\|_{\infty}$ 是作为连续函数空间的范数.此时令
\begin{equation*}
\begin{aligned}
&\|\varphi\|_{V} = \inf\left\{\sum_{j=1}^{\infty}\|f_{j}\|_{\infty}\|g_{j}\|_{\infty}: \varphi(x, y) = \sum_{j=1}^{\infty}f_{j}(x)g_{j}(y),\right.\\
&\left.\sum_{j=1}^{\infty}\|f_{j}\|_{\infty}\|g_{j}\|_{\infty} < \infty\right\}.
\end{aligned}
\end{equation*}
$V(K_{1} \times K_{2}) = C(K_{1}) \widehat{\otimes}\, C(K_{2})$ 也记为 $V(K_{1} \times K_{2})$.
\end{definition}

\begin{proposition}\label{prop:3.4.7}
$V(K_{1} \times K_{2})$ 关于 $\|\cdot\|_{V}$ 是半单,正则的交换 Banach 代数,其极大理想空间为 $K_{1} \times K_{2}$.
\end{proposition}

\begin{proof}
$V(K_{1} \times K_{2})$ 关于 $\|\cdot\|_{V}$ 成为交换 Banach 代数这一点是容易证明的.又,若 $(x, y) \in K_{1} \times K_{2}$,则由 $V(K_{1} \times K_{2}) \ni \varphi \mapsto \varphi(x, y) \in \mathbf{C}$ 可把 $K_{1} \times K_{2}$ 视为 $\Delta(V(K_{1} \times K_{2}))$ 的子集.反之,若 $h \in \Delta(V(K_{1} \times K_{2}))$,则存在 $\varphi_{0} \in V(K_{1} \times K_{2})$ 使得 $h(\varphi_{0}) \neq 0$.此时,对 $f \in C(K_{1})$,$g \in C(K_{2})$ 令
\begin{equation*}
h_{1}(f) = \frac{h(f\varphi_{0})}{h(\varphi_{0})},\quad h_{2}(g) = \frac{h(g\varphi_{0})}{h(\varphi_{0})},
\end{equation*}
则 $h_{1} \in \Delta(C(K_{1}))$,$h_{2} \in \Delta(C(K_{2}))$.事实上,若 $h(\varphi') \neq 0$,则由
\begin{equation*}
\frac{h(f\varphi_{0})}{h(\varphi_{0})} = \frac{h(f\varphi_{0})h(\varphi')}{h(\varphi_{0}\varphi')} = \frac{h(f\varphi_{0}\varphi')}{h(\varphi_{0}\varphi')} = \frac{h(f\varphi')}{h(\varphi')}
\end{equation*}
成立,可得
\begin{align*}
h_{1}(f_{1}f_{2}) &= \frac{h(f_{1}f_{2}\varphi_{0}^{2})}{h(\varphi_{0}^{2})} = \frac{h(f_{1}\varphi_{0})h(f_{2}\varphi_{0})}{h(\varphi_{0})h(\varphi_{0})} = h_{1}(f_{1})h_{1}(f_{2})\\
&\quad (f_{1}, f_{2} \in C(K_{1})).
\end{align*}
因此,由 $\mathfrak{M}(C(K_{1})) = K_{1}$,存在 $x_{1} \in K_{1}$ 使得 $h_{1}(f) = f(x_{1})$($f \in C(K_{1})$).对 $h_{2}$ 也同样.于是
\begin{equation*}
\begin{aligned}
i: K_{1} \times K_{2} &\ni (x_{1}, x_{2}) \to h_{(x_{1}, x_{2})} \in \Delta(V(K_{1} \times K_{2}))\\
&(h_{(x_{1}, x_{2})}(\varphi) = \varphi(x_{1}, x_{2}) \in V(K_{1} \times K_{2}))
\end{aligned}
\end{equation*}
是双射.而且由于它连续,故是同胚映射,$K_{1} \times K_{2}$ 与 $\mathfrak{M}(V(K_{1} \times K_{2}))$ 在自然对应下等同.此外,半单,正则也是容易看出的.
\end{proof}

这里令 $V(\mathbf{T}^{3}) = V(\mathbf{T}^{3} \times \mathbf{T}^{3})$,关于 $V(\mathbf{T}^{3})$ 的谱综合,下式成立.

\begin{theorem}\label{thm:3.4.8}
$V(\mathbf{T}^{3})$ 不可谱综合.
\end{theorem}

\begin{proof}
从映射 $M$,$P$ 的定义开始.如下定义 $M: C(\mathbf{T}^{3}) \to C(\mathbf{T}^{3} \times \mathbf{T}^{3})$,$P: C(\mathbf{T}^{3} \times \mathbf{T}^{3}) \to C(\mathbf{T}^{3})$.其中 $dm(y)$ 是 $\mathbf{T}^{3}$ 上的 Lebesgue 测度,且 $m(\mathbf{T}^{3}) = 1$.
\begin{equation*}
(Mf)(x, y) = f(x + y), \quad (P\varphi)(x) = \int_{\mathbf{T}^{3}}\varphi(x - y, y)\dd m(y).
\end{equation*}
此时,若 $f \in C(\mathbf{T}^{3})$ 则 $Mf(x, y) \in C(\mathbf{T}^{3} \times \mathbf{T}^{3})$;若 $\varphi \in C(\mathbf{T}^{3} \times \mathbf{T}^{3})$ 则 $P\varphi \in C(\mathbf{T}^{3})$.又由
\begin{equation*}
(PMf)(x) = \int_{\mathbf{T}^{3}}(Mf)(x - y, y)\dd m(y) = \int_{\mathbf{T}^{3}}f(x)\dd m(y) = f(x)
\end{equation*}
$PM$ 是恒等映射.利用这一点来证明 $V(\mathbf{T}^{3})$ 不可谱综合.

由定理~\ref{thm:3.4.5},$A(\mathbf{T}^{3})$ 不可谱综合,故存在非谱综合集 $E$ 和 $f \in A(\mathbf{T}^{3})$ 使得 $f(x) = 0$($x \in E$),且 $f$ 不能被在 $E$ 的邻域上为零的 $g \in A(\mathbf{T}^{3})$ 所逼近.这里令 $E^{*} = \{(x, y): x + y \in E\}$,由 $f \in A(\mathbf{T}^{3})$ 可写成 $f(x) = \sum_{n \in \mathbf{Z}^{3}} a_{n}e^{inx}$,$\sum_{n \in \mathbf{Z}^{3}}|a_{n}| < \infty$.此时
\begin{equation*}
Mf(x, y) = f(x + y) = \sum_{n \in \mathbf{Z}^{3}}a_{n}e^{inx}e^{iny} = \sum_{n \in \mathbf{Z}^{3}}(a_{n}e^{inx})e^{iny}
\end{equation*}
所以 $Mf \in V(\mathbf{T}^{3})$.而且 $Mf(z) = 0$($z \in E^{*}$).现在,设 $\psi(x, y) = \sum f_{j}(x)g_{j}(y) \in V(\mathbf{T}^{3})$ 在 $E^{*}$ 的邻域上为零,则
\begin{equation*}
P\psi(x) = \int \psi(x - y, y)\dd m(y) = \sum_{j} f_{j} * g_{j}(x),
\end{equation*}
即得 $P\psi(x) = \sum f_{j} * g_{j}(x)$.这里由 $f_{j}, g_{j} \in L^{2}(\mathbf{T}^{3})$,
\begin{equation*}
\|f_{j} * g_{j}\|_{A(\mathbf{T}^{3})} \leq \|\widehat{f_{j}}\|_{2}\|\widehat{g_{j}}\|_{2} \leq \|f_{j}\|_{\infty}\|g_{j}\|_{\infty},
\end{equation*}
且 $\sum\|f_{j} * g_{j}\|_{A(\mathbf{T}^{3})} < \infty$,故 $P\psi \in A(\mathbf{T}^{3})$,且 $\sup_{\|\psi\|_{V} \leq 1}\|P\psi\|_{A(\mathbf{T}^{3})} \leq 1$.因此
\begin{equation*}
\|f - P\psi\|_{A(\mathbf{T}^{3})} = \|P(Mf - \psi)\|_{A(\mathbf{T}^{3})} \leq \|Mf - \psi\|_{V(\mathbf{T}^{3})}.
\end{equation*}
又,$P\psi$ 在 $E$ 的邻域上为 0.从而 $\inf\{\|Mf - \psi\|_{V(\mathbf{T}^{3})}: \psi \in V(\mathbf{T}^{3}),\ \psi \text{ 在 } E^{*} \text{ 的邻域上为零}\} > 0$,由 $E$ 的性质,$V(\mathbf{T}^{3})$ 不可谱综合.
\end{proof}

\subsection{\texorpdfstring{$A(\mathbf{T})$ 的谱综合}{A(T) 的谱综合}}\label{sec:3.4.4}

这里,用 Varopoulos 的方法 \cite{ch3:52}(cf. \cite{ch3:23})证明 $A(\mathbf{T})$ 不可谱综合.如后所述,该方法并不限于 $\mathbf{T}$,对非离散紧 Abel 群也可作类似的讨论.首先,设 $D_{2}$ 为循环群 $\{0, 1\}$ 的可数直积所构成的紧 Abel 群.把 $D_{2}$ 表示为 $D_{2} = \prod_{j=1}^{\infty}\{0, 1\}$,拓扑取直积拓扑.此时,因 $D_{2}$ 是 Cantor 集,利用这一点如下推进.如 3.3.3 小节所述,$\mathbf{T}$ 中存在作为 Cantor 集的 Kronecker 集,记之为 $K$.把 $K$ 分解为 $K = K_{1} \cup K_{2}$,$K_{1} \cap K_{2} = \emptyset$,其中 $K_{1}, K_{2}$ 是作为 Cantor 集的 Kronecker 集.本小节参考了 \cite{ch3:23}.

\begin{definition}\label{def:3.4.9}
如下定义:
(1) $\widetilde{K} = K_{1} + K_{2} = \{k_{1} + k_{2}: k_{j} \in K_{j},\ j = 1, 2\}$.此时,由 Kronecker 集的线性无关性,$k \in K$ 可唯一表示为 $k = k_{1} + k_{2}$($k_{j} \in K_{j}$,$j = 1, 2$).

(2) 注意:因 $D_{2}$ 是 Cantor 集,$D_{2}$ 与 $K_{j}$($j = 1, 2$)可等同,利用这一点可把 $D_{2} \times D_{2}$ 视为 $K_{1} + K_{2} = \widetilde{K}$.

(3) 当 $E \subset \mathbf{T}$ 为紧集时,把 $A(\mathbf{T})$ 到 $E$ 的限制 $A(E)$ 定义为商环 $A(E) = A(\mathbf{T})/I(E)$.
\end{definition}

\begin{proposition}\label{prop:3.4.10}
若 $\mathbf{T}$ 的紧集 $E$ 是 Kronecker 集,则 $A(E) = A(\mathbf{T})/I(E)$ 与 $E$ 上的连续函数空间 $C(E)$ 作为 Banach 空间同构.
\end{proposition}

\begin{proof}
把 $A(E)$ 看作 $A(\mathbf{T})$ 中函数在 $E$ 上的限制,只需证明 $A(E) = C(E)$.为证明这一点,由 Riesz 定理 $C(E)^{*} = M(E)$,只需对 $\mu \in M(E)$ 证明
\begin{equation*}
\|\mu\|_{M(E)} \leq C\|\widehat{\mu}\|_{\infty}.
\end{equation*}
不过这里令 $\|\widehat{\mu}\|_{\infty} = \sup_{m \in \mathbf{Z}}|\widehat{\mu}(m)|$.若这已证明,则因 $A(E)$ 与 $C(E)$ 都是 Banach 空间,由泛函分析中的值域定理可得 $A(E) = C(E)$.于是,对 $\mu \in M(E)$,取适当的可测函数 $f(x)$ 使 $|f(x)| = 1$($x \in E$),有 $\|\mu\| = \left|\int f(x)\dd\mu(x)\right|$;而由 $\mu \in M(E)$,$f(x)$ 可取为连续函数.又因 $E$ 是 Kronecker 集,$f(x)$ 可由适当的 $e^{imx}$ 在 $E$ 上一致逼近,故 $\|\mu\| \leq \|\widehat{\mu}\|_{\infty}$ 成立.
\end{proof}

\begin{proposition}\label{prop:3.4.11}
$A(\widetilde{K})$ 作为 Banach 代数与 $V(D_{2}) = C(D_{2}) \widehat{\otimes}\, C(D_{2})$ 同构.
\end{proposition}

\begin{proof}
注意:由定义,$\varphi \in V(D_{2})$ 可表示为 $D_{2} \times D_{2}$ 上的连续函数.

首先证明 $A(\widetilde{K}) \subset V(D_{2})$.$f \in A(\widetilde{K})$ 可写成 $f(z) = \sum_{j \in \mathbf{Z}} a_{j}e^{ijz}$($z \in \widetilde{K}$),$\sum_{j \in \mathbf{Z}}|a_{j}| < \infty$.此时,若设 $z = x + y$,$x \in K_{1}$,$y \in K_{2}$,则
\begin{equation*}
\varphi(x, y) = f(x + y) = \sum_{j \in \mathbf{Z}}a_{j}e^{ijx}e^{ijy}, \quad \sum_{j \in \mathbf{Z}}\|a_{j}e^{ijx}\|_{\infty}\|e^{ijy}\|_{\infty} < \infty,
\end{equation*}
所以 $\varphi(x, y) \in V(K_{1} \times K_{2}) = C(K_{1}) \widehat{\otimes}\, C(K_{2}) = C(D_{2}) \widehat{\otimes}\, C(D_{2}) = V(D_{2})$.于是证明了 $A(\widetilde{K}) \subset V(D_{2})$.

反之,若 $\varphi \in V(D_{2})$,则 $\varphi(x, y) = \sum_{j=1}^{\infty}f_{j}(x)g_{j}(y)$,$\sum_{j}\|f_{j}\|_{\infty}\|g_{j}\|_{\infty} \leq 2\|\varphi\|_{V(D_{2})}$.此时把 $K_{1} = D_{1}, K_{2} = D_{2}$ 看待,则 $\varphi(x, y) \in C(K_{1} \times K_{2})$.又因 $K_{1}, K_{2}$ 是 Kronecker 集,由命题~\ref{prop:3.4.10} 的证明,存在常数 $C > 0$ 使得 $\|f_{j}\|_{A(K_{1})} \leq C\|f_{j}\|_{C(K_{1})}$,$\|g_{j}\|_{A(K_{2})} \leq C\|g_{j}\|_{C(K_{2})}$.这里令
\begin{equation*}
f_{j}(x) = \sum_{l \in \mathbf{Z}}a_{jl}e^{ilx}, \quad \sum_{l}|a_{jl}| \leq \left(\frac{3}{2}\right)\|f_{j}\|_{A(K_{1})},
\end{equation*}
\begin{equation*}
g_{j}(x) = \sum_{m \in \mathbf{Z}}b_{jm}e^{imx}, \quad \sum_{m}|b_{jm}| \leq \left(\frac{3}{2}\right)\|g_{j}\|_{A(K_{2})}
\end{equation*}
($j = 1, 2, \ldots$),则
\begin{align*}
\sum_{j,l,m}|a_{jl}||b_{jm}| &\leq \frac{9}{4}\sum_{j}\|f_{j}\|_{A(K_{1})}\|g_{j}\|_{A(K_{2})}\\
&\leq \frac{9}{4}C^{2}\sum_{j}\|f_{j}\|_{\infty}\|g_{j}\|_{\infty} \leq 5C^{2}\|\varphi\|_{V(D_{2})}.
\end{align*}
进一步,令
\begin{equation*}
f(z) = \left\{\begin{array}{ll}
e^{ilz} & (z \in K_{1})\\
e^{imz} & (z \in K_{2})
\end{array}\right.
\end{equation*}
因 $K = K_{1} \cup K_{2}$ 是 Kronecker 集,对 $\varepsilon > 0$ 存在 $k_{lm} \in \mathbf{Z}$ 使得 $|f(z) - e^{ik_{lm}z}| < \varepsilon$($z \in K$).现在令
\begin{equation*}
T_{\varepsilon}\varphi(z) = \sum_{j,l,m}a_{jl}b_{jm}e^{ik_{lm}z},
\end{equation*}
则
\begin{equation*}
\|T_{\varepsilon}\varphi\|_{A(\widetilde{K})} \leq \sum_{j,l,m}|a_{jl}||b_{jm}| \leq 5C^{2}\|\varphi\|_{V(D_{2})}.
\end{equation*}
又由 $|e^{ilx}e^{imy} - e^{ik_{lm}(x+y)}| < 2\varepsilon$($x \in K_{1}$,$y \in K_{2}$),
\begin{equation*}
\|T_{\varepsilon}\varphi - \varphi\|_{V(D_{2})} \leq 2\varepsilon\sum_{j,l,m}|a_{jl}||b_{jm}| \leq 10C^{2}\varepsilon\|\varphi\|_{V(D_{2})}.
\end{equation*}
此时,考虑级数
\begin{equation}\tag{3.4.1}
T_{\varepsilon}\varphi + T_{\varepsilon}(I - T_{\varepsilon})\varphi + \cdots + T_{\varepsilon}(I - T_{\varepsilon})^{m}\varphi + \cdots
\label{eq:3.4.1}\end{equation}
其中 $I$ 是恒等映射.在这种情况下,因
\begin{equation*}
\|T_{\varepsilon}(I - T_{\varepsilon})^{m}\varphi\|_{A(\widetilde{K})} \leq 5C^{2}\|(I - T_{\varepsilon})^{m}\varphi\|_{V(D_{2})} \leq 5C^{2}(10C^{2}\varepsilon)^{m}\|\varphi\|_{V(D_{2})},
\end{equation*}
当 $\varepsilon$ 充分小时\eqref{eq:3.4.1}在 $A(\widetilde{K})$ 中收敛.又由 $\varphi - T_{\varepsilon}\varphi = \sum_{m=1}^{N}T_{\varepsilon}(I - T_{\varepsilon})^{m}\varphi + (I - T_{\varepsilon})^{N+1}\varphi$,$\|(I - T_{\varepsilon})^{N+1}\varphi\|_{V(D_{2})} \leq (10C^{2}\varepsilon)^{N+1}\|\varphi\|_{V(D_{2})} \to 0$($N \to \infty$),得 $\varphi = \sum_{m=0}^{\infty}T_{\varepsilon}(I - T_{\varepsilon})^{m}\varphi$,所以 $\varphi \in A(\widetilde{K})$.
\end{proof}

\begin{theorem}\label{thm:3.4.12}
$A(\widetilde{K})$ 不可谱综合.
\end{theorem}

\begin{proof}
设 $T = \{e^{i2\pi x}: 0 \leq x < 1\}$,如下定义从 $D_{2}$ 到 $T^{3}$ 的映射 $\tau$:
\begin{equation*}
\tau: (\varepsilon_{1}, \varepsilon_{2}, \ldots) \to (e^{2\pi i(0.\varepsilon_{1}\varepsilon_{4}\cdots)}, e^{2\pi i(0.\varepsilon_{2}\varepsilon_{5}\cdots)}, e^{2\pi i(0.\varepsilon_{3}\varepsilon_{6}\cdots)}).
\end{equation*}
这里 $\varepsilon_{j} = 0, 1$,$0.\varepsilon_{j_{1}}\varepsilon_{j_{2}}\cdots$ 表示由二进制展开得到的区间 $[0, 1)$ 的元素.这一从 $D_{2}$ 到 $T^{3}$ 的对应连续,且除 $D_{2}$ 的可数个集合外是一一的.由此在函数空间之间诱导的映射一律用 $\widetilde{\tau}$(即 $\widetilde{\tau}F = F(\tau)$)表示.又由 $\tau$ 的对应,$C(D_{2})$ 自然被视为 $L^{\infty}(\mathbf{T}^{3})$ 的子集.把这个嵌入记为 $\widetilde{i}$.另外,由 $\tau$ 诱导的从 $V(D_{2})$ 到 $L^{\infty}(\mathbf{T}^{3}) \widehat{\otimes}\, L^{\infty}(\mathbf{T}^{3})$ 的映射也用 $\widetilde{i}$ 表示.进一步如下定义交换 Banach 代数 $L^{\infty}(\mathbf{T}^{3}) \widehat{\otimes}\, L^{\infty}(\mathbf{T}^{3})$:称 $\phi \in L^{\infty}(\mathbf{T}^{3}) \widehat{\otimes}\, L^{\infty}(\mathbf{T}^{3})$,是指
\begin{align*}
\phi(x, y) &= \sum_{j=1}^{\infty}f_{j}(x)g_{j}(y) \quad (f_{j}(x), g_{j}(x) \in L^{\infty}(\mathbf{T}^{3})),\\
&\sum_{j}\|f_{j}\|_{L^{\infty}}\|g_{j}\|_{L^{\infty}} < \infty.
\end{align*}
此时,由
\begin{equation*}
\|\phi\|_{L^{\infty}(\mathbf{T}^{3}) \widehat{\otimes} L^{\infty}(\mathbf{T}^{3})} = \inf\left\{\sum_{j=1}^{\infty}\|f_{j}\|_{L^{\infty}}\|g_{j}\|_{L^{\infty}}: \text{上述表示}\right\}
\end{equation*}
$L^{\infty}(\mathbf{T}^{3}) \widehat{\otimes}\, L^{\infty}(\mathbf{T}^{3})$ 如 3.4.3 小节那样成为交换 Banach 代数.

于是,利用 $\widetilde{\tau}, \widetilde{i}$,
\begin{align*}
V(\mathbf{T}^{3}) = C(\mathbf{T}^{3}) \widehat{\otimes}\, C(\mathbf{T}^{3}) &\xrightarrow{\widetilde{\tau}} C(D_{2}) \widehat{\otimes}\, C(D_{2}) = V(D_{2})\\
&\xrightarrow{\widetilde{i}} L^{\infty}(\mathbf{T}^{3}) \widehat{\otimes}\, L^{\infty}(\mathbf{T}^{3})
\end{align*}
中各算子作用素范数不超过 1.此时,由定理~\ref{thm:3.4.8},$V(\mathbf{T}^{3})$ 不可谱综合.设闭集 $E \subset \mathbf{T}^{3} \times \mathbf{T}^{3}$ 是 $V(\mathbf{T}^{3})$ 的非谱综合集,则存在在 $E$ 上为零,且不能被在 $E$ 的邻域上为零的 $V(\mathbf{T}^{3})$ 的元逼近的 $f \in V(\mathbf{T}^{3})$.此时,$(\tau \times \tau)^{-1}(E)$ 是 $D_{2} \times D_{2}$ 的闭集,由 $\widetilde{\tau}f \in V(D_{2})$ 及 $f$ 的取法,$\widetilde{\tau}f(z) = 0$($z \in (\tau \times \tau)^{-1}(E)$).若设 $h \in V(D_{2})$ 在 $(\tau \times \tau)^{-1}(E)$ 的邻域上为零,且逼近了 $\widetilde{\tau}f$.此时,令 $\tilde{i}h(x, y) = \sum_{j}h_{1j}(x)h_{2j}(y)$,则由 $\tilde{i}h(x, y) = \sum(\tilde{i}h_{1j})(x)(\tilde{i}h_{2j})(y)$,$\tilde{i}h(x, y)$ 可视为 $L^{\infty}(\mathbf{T}^{3}) \widehat{\otimes}\, L^{\infty}(\mathbf{T}^{3})$ 的元.于是,定义函数列 $\{e_{n}(x, y)\}_{n}$ 为 $e_{n}(x, y) = e_{1n}(x)e_{2n}(y)$.其中,函数列 $\{e_{jn}\}_{n}$($j = 1, 2$)满足 $e_{jn} \in C(\mathbf{T}^{3})$,$e_{jn} \geq 0$,$\{\supp e_{jn}\}_{n}$ 是递减列,其交集为 $\{0\}$,而且 $\int e_{jn}(x)\dd x = 1$.这里,考虑
\begin{equation*}
(e_{n} * \widetilde{i}h)(x, y) = \sum_{j}(e_{1n} * (\widetilde{i}h_{1j}))(x)(e_{2n} * (\widetilde{i}h_{2j}))(y),
\end{equation*}
则 $e_{n} * \widetilde{i}h$ 在 $E$ 的邻域上为零,而且由构造方法,$e_{n} * \widetilde{i}h \in V(\mathbf{T}^{3})$,$\|e_{n} * \widetilde{i}h - f\|_{V(\mathbf{T}^{3})} \to 0$($n \to \infty$).这与 $E$ 的假设矛盾.因此 $(\tau \times \tau)^{-1}(E)$ 是 $V(D_{2})$ 的非谱综合集,$V(D_{2})$ 不可谱综合.从而由命题~\ref{prop:3.4.11},$A(\widetilde{K})$ 也不可谱综合.
\end{proof}

\begin{theorem}\label{thm:3.4.13}
$A(\mathbf{T})$ 不可谱综合.
\end{theorem}

\begin{proof}
由定理~\ref{thm:3.4.12},$A(\widetilde{K})$ 不可谱综合.因此,存在 $A(\widetilde{K})$ 的非谱综合集 $E_{0} \subset \widetilde{K} = K_{1} + K_{2}$.此时,假设 $E_{0}$ 是 $A(\mathbf{T})$ 的谱综合集.若 $f(x) = 0$($x \in E_{0}$),$f \in A(\widetilde{K})$,则存在 $\widetilde{f} \in A(\mathbf{T})$ 使得 $\widetilde{f}(x) = f(x)$($x \in \widetilde{K}$).由 $\widetilde{f}(x) = 0$($x \in E_{0}$)且 $E_{0}$ 是 $A(\mathbf{T})$ 的谱综合集,对 $\varepsilon > 0$ 存在 $\widetilde{g} \in A(\mathbf{T})$,$\widetilde{g}$ 在 $E_{0}$ 的 $\mathbf{T}$ 中的邻域上为零,且 $\|\widetilde{f} - \widetilde{g}\|_{A(\mathbf{T})} < \varepsilon$.因此,若令 $g(x) = \widetilde{g}(x)$($x \in \widetilde{K}$),则 $g \in A(\widetilde{K})$,$g$ 在 $\widetilde{K}$ 的相对拓扑下于 $E_{0}$ 的邻域上为零,且 $\|f - g\|_{A(\widetilde{K})} < \varepsilon$,这与 $E_{0}$ 是 $A(\widetilde{K})$ 的非谱综合集矛盾.故 $E_{0}$ 是 $A(\mathbf{T})$ 的非谱综合集,从而证明了 $A(\mathbf{T})$ 不可谱综合.
\end{proof}

\begin{theorem}\label{thm:3.4.14}
$A(\mathbf{R})$ 不可谱综合.
\end{theorem}

\begin{proof}
设 $E_{0} \subset T$ 为定理~\ref{thm:3.4.13} 中的 $A(\mathbf{T})$ 的非谱综合集,$E_{0}, \widetilde{K}$ 可取为 $\{|x| < \frac{\pi}{2}\}$ 的子集.假设 $E_{0}$ 是 $A(\mathbf{R})$ 的谱综合集.此时,设 $f \in A(\widetilde{K})$,$f(x) = 0$($x \in E_{0}$).$f$ 可用同一记号看作 $A(\mathbf{T})$ 到 $\widetilde{K}$ 的限制.现在,取 $h \in A(\mathbf{R})$ 使得 $h$ 在 $\widetilde{K}$ 的邻域上为 1 且 $\supp h \subset \{|x| < \frac{\pi}{2}\}$.此时,由命题~\ref{prop:3.4.4} 的想法,$fh \in A(\mathbf{R})$.因 $(fh)(x) = 0$($x \in E_{0}$),由 $E_{0}$ 是 $A(\mathbf{R})$ 的谱综合集,对 $\varepsilon > 0$ 存在 $\widetilde{g} \in A(\mathbf{R})$,$\widetilde{g}$ 在 $E_{0}$ 的邻域上为零且 $\|fh - \widetilde{g}\|_{A(\mathbf{R})} < \varepsilon$.于是取满足 $k(x) = 1$($x \in \supp h$),$\supp k \subset \{|x| < \frac{\pi}{2}\}$,$\|k\|_{A(\mathbf{R})} \leq 2$ 的 $k$,则由命题~\ref{prop:3.4.4} 的想法,$k\widetilde{g}$ 可视为 $A(\mathbf{T})$ 的元,故
\begin{align*}
\|k(\widetilde{g} - fh)\|_{A(\widetilde{K})} &\leq \|k(\widetilde{g} - fh)\|_{A(\mathbf{T})}\\
&\leq C\|k(\widetilde{g} - fh)\|_{A(\mathbf{R})} \leq 2C\|\widetilde{g} - fh\|_{A(\mathbf{R})} < 2C\varepsilon.
\end{align*}
其中 $C$ 是适当的常数.这里 $(k\widetilde{g} - fh)(x) = \widetilde{g}(x) - f(x)$($x \in \widetilde{K}$),且 $\widetilde{g}(x)$ 在 $E_{0}$ 的邻域上为零.综上,$\|\widetilde{g} - f\|_{A(\widetilde{K})} < 2C\varepsilon$,$\widetilde{g}$ 在 $E_{0}$ 的邻域上为零,与 $E_{0}$ 是 $A(\widetilde{K})$ 的非谱综合集矛盾.因此 $E_{0}$ 是 $A(\mathbf{R})$ 的非谱综合集.
\end{proof}

\begin{note}\label{note:3.4.15}
$A(\mathbf{T})$,$A(\mathbf{R})$ 不可谱综合的 Varopoulos 证明方法,依赖于作为 Cantor 集的 Kronecker 集的存在.因此,对于 $\mathbf{T}$ 以外的,可用这一方法进行讨论的非离散紧 Abel 群,也可作类似的讨论(cf. \cite{ch3:52}).
\end{note}

\subsection{\texorpdfstring{$A(\mathbf{T})$ 的非谱综合集}{A(T) 的非谱综合集}}\label{sec:3.4.5}

\begin{definition}\label{def:3.4.16}
紧集 $E \subset \mathbf{T}$ 称为 Helson 集,是指把 $A(\mathbf{T})/I(E)$ 看作 $A(\mathbf{T})$ 到 $E$ 的限制 $A(E)$ 时,$A(E) = C(E)$ 成立.
\end{definition}

即,$E$ 上任意连续函数都是具有绝对收敛 Fourier 级数的连续函数在 $E$ 上的限制.

从而,由命题~\ref{prop:3.4.10},Kronecker 集是 Helson 集.

如 3.3.3 小节所述,Varopoulos \cite{ch3:50} 证明了 Kronecker 集是 $A(\mathbf{T})$ 的谱综合集.其次,Helson 集的情形成为问题.Körner \cite{ch3:33} 于 1972 年,对长期未解决的``Helson 集是否为谱综合集''的问题,证明了``存在既是 Helson 集又是非谱综合集的集合'',给予了否定的解答.其后不久,R. Kaufman 对 T. W. Körner 的证明给出了简洁的别证.进而,S. Saeki 推广了该证明.这里只介绍结果.详情请参阅 R. Kaufman 的论文或 Saeki \cite{ch3:45}(cf. \cite{ch3:46}).

下面是由 Rudin \cite{ch3:11} 给出的(cf. \cite{ch3:33}, \cite{ch3:42}).

\begin{proposition}\label{prop:3.4.17}
$\mathbf{T}$ 存在具有如下性质的完全不连通的紧子集 $E$:

(1) $E$ 的 Lebesgue 测度为零.

(2) $\{\sigma \in M(E): \lim_{|n| \to \infty}\widehat{\sigma}(n) = 0,\ \sigma \neq 0\} \neq \emptyset$.
\end{proposition}

关于这一点,Kaufman \cite{ch3:32}(cf. \cite{ch3:45})证明了下面的事实.

\begin{theorem}[{$\;$\cite{ch3:32}(cf. \cite{ch3:45})}]\label{thm:3.4.18}
在命题~\ref{prop:3.4.17} 的记号下,存在满足 $E_{0} \subset \operatorname{supp}\sigma$ 的 Helson 集 $E_{0}$,使得 $\{T \in PM(E_{0}): \lim_{|n| \to \infty}\widehat{T}(n) = 0,\ T \neq 0\} \neq \emptyset$.其中,设 $PM(E_{0}) = \{s \in PM(\mathbf{T}): \supp s \subset E_{0}\}$,$\supp s$ 是指满足下列条件的开集 $O$ 的补集的最小者:对所有支集含于 $O$ 的 $f \in A(\mathbf{T})$,$(s, f) = 0$.
\end{theorem}

\begin{corollary}[{$\;$\cite{ch3:32},cf. \cite{ch3:45}}]\label{cor:3.4.19}
$\mathbf{T}$ 中存在既是 Helson 集又是非谱综合集的集合.
\end{corollary}

由定理~\ref{thm:3.4.18},系~\ref{cor:3.4.19} 的证明如下进行.

\begin{proof}[系~\ref{cor:3.4.19} 的证明]
设 $E_{0}$ 为定理~\ref{thm:3.4.18} 中的 Helson 集.此时,作为 Helson 集的性质,已知
\begin{equation*}
\left\{\mu \in M(\mathbf{T}): \lim_{|n| \to \infty}\widehat{\mu}(n) = 0,\ \mu \neq 0\right\} = \emptyset
\end{equation*}
成立(cf. \cite{ch3:11}).若 $E_{0}$ 是 $A(\mathbf{T})$ 的谱综合集,则定理~\ref{thm:3.4.18} 中的 $T$ 是 $(A(\mathbf{T})/J(E_{0}))^{*}$ 的元,由 $I(E_{0}) = J(E_{0})$ 得 $T \in A(E_{0})^{*}$.因 $E_{0}$ 是 Helson 集,$A(E_{0}) = C(E_{0})$,由 Riesz 表示定理得 $T \in M(E_{0})$.这与证明开头所述相矛盾.因此 $E_{0}$ 是非谱综合集.
\end{proof}

\subsection{Schwartz 结果的发展}\label{sec:3.4.6}

如前所述,L. Schwartz 证明了 $S^{n-1}$($n \geq 3$)是交换 Banach 代数 $A(\mathbf{R}^{n})$ 中的非谱综合集.将其发展的研究见于 N. Th. Varopoulos,Y. Domar,D. Müller 等人.这里简单介绍 Müller \cite{ch3:35}.详情请连同文献参阅 \cite{ch3:35}.以下使用 3.1 节的记号.

\begin{definition}\label{def:3.4.20}
(1) 设 $M$ 是 $\mathbf{R}^{n}$ 的 $k$ 维子流形,对 $\alpha \in M$ 设 $(X, \Omega)$ 为坐标邻域.即 $\Omega$ 是 $\mathbf{R}^{k}$ 的开集,$X \in C^{\infty}(\Omega, \mathbf{R}^{n})$,$X(x_{0}) = \alpha$($0 \leq k \leq n - 1$).此时,设 $N$ 为 $\alpha$ 处的法向量,令
\begin{equation*}
h_{ij}(\alpha; X, N) = N \cdot \frac{\partial^{2}}{\partial x_{i}\partial x_{j}}X(x_{1}, \ldots, x_{k}) = \frac{\partial^{2}}{\partial x_{i}\partial x_{j}}(N \cdot X),
\end{equation*}
由此定义第二基本矩阵 $(h_{ij}(\alpha; X, N))_{i,j}$.

(2) 令 $\pi(\alpha; X, N) = \{y \in \mathbf{R}^{k}: \sum_{j=1}^{k}h_{ij}y^{j} = 0,\ i = 1, \ldots, k\}$,$\pi(y) = (h_{ij})(y_{i})$,$y = (y_{i})$,则 $\pi$ 可看作从 $\alpha$ 处的切平面 $M_{\alpha}$ 到自身的映射,$\pi(\alpha; X, N)$ 可看作切平面 $M_{\alpha}$ 的子集.而且 $\dim\pi(\alpha; X, N) = k - \operatorname{rank}(h_{ij})$.另请注意 $\pi(\alpha; X, N)$ 不依赖于 $X$ 的取法.$y \in \pi(\alpha; X, N)$ 意味着:对过 $\alpha$ 的 $M$ 上的曲线 $p(t)$($p(0) = \alpha$,$\dot{p}(0) = y$,$-1 \leq t \leq 1$),沿 $p(t)$ 的法向量 $\{N(x)\}$ 不``摆动''.

(3) 令
\begin{equation*}
\pi(\alpha; X) = \bigcap_{N: \alpha \text{ 处的法向量}}\pi(\alpha; X, N).
\end{equation*}
由 (2) 的注意,$\pi(\alpha; X)$ 不依赖于 $X$ 的取法,所以 $\dim(\pi(\alpha; X))$ 也不依赖于坐标邻域.称 $\dim\pi(\alpha; X)$ 为 $M$ 在 $\alpha$ 处的 relative nullity(相对零化度),记为 $\nu(\alpha) = \dim\pi(\alpha; X)$.
\end{definition}

\begin{example}\label{exa:3.4.21}
当 $n = 3$ 时,即在 $\mathbf{R}^{3}$ 中如下:

(i) 当 $M = S^{2}$ 时,由 $\pi(\alpha; X, N) = \{0\}$ 得 $\nu(\alpha) = 0$($\alpha \in M$).

(ii) 当 $M = S^{1}$(平面的单位圆)时,$\nu(\alpha) = 0$($\alpha \in M$).

(iii) 当 $M$ 为圆柱面时,$\nu(\alpha) = 1$($\alpha \in M$).
\end{example}

\begin{definition}\label{def:3.4.22}
设 $M$ 是 $\mathbf{R}^{n}$ 的 $k$ 维流形,对 $\alpha \in M$ 设 $(X, \Omega)$ 为 $\alpha$ 处的坐标邻域系,$N_{0}$ 为 $\alpha$ 处的法向量.此时如下定义 $\alpha$ 处的曲率 $K(\alpha)$:
\begin{equation*}
K(\alpha) = \frac{\det(h_{ij}(\alpha; X, N_{0}))}{\det\left(\frac{\partial X}{\partial x_{i}}, \frac{\partial X}{\partial x_{j}}\right)}.
\end{equation*}
这里注意:因 $M$ 是 $k$ 维流形,分母不为零.又,$K(\alpha) = 0$ 的充要条件是分子为零,即 $\nu(\alpha) = k - \operatorname{rank}(h_{ij}) > 0$.
\end{definition}

以下设 $n = 3$,在 $\mathbf{R}^{3}$ 的空间中讨论.

\begin{definition}\label{def:3.4.23}
$F \subset \mathbf{R}^{2}$ 在 $y_{0} \in \mathbf{R}^{2}$ 具有 restricted cone property(限制锥性质),是指:对 $y_{0}$ 的邻域 $V_{0}$,某个 $0 < \sigma < 1$,$y_{1} \in \mathbf{R}^{2}$($\|y_{1}\| = 1$),令 $K = \{y \in \mathbf{R}^{2}: (1 - \sigma)\|y\| \leq y \cdot y_{1} \leq \sigma\}$ 时,$F \cap V_{0} - K \subset F$ 成立.其中 $\|z\|$ 表示 $z$ 与原点的距离.
\end{definition}

\begin{definition}\label{def:3.4.24}
(1) 紧集 $E(\subset M)$ 具有 restricted cone property(限制锥性质),是指:对任意 $\alpha \in E$,存在 $\alpha$ 的邻域 $V$,使得 $E \cap V$ 向 $\alpha$ 处切平面的投影在 $\alpha$ 具有 restricted cone property(限制锥性质).

(2) 令 $J_{1}(E) = \{\varphi \in C_{c}^{\infty}(\mathbf{R}^{3}): \varphi(x) = 0 \ (x \in E)\}$,对 $l = 1, 2, \ldots$,把 $J_{1}(E)^{l}$ 定义为由 $J_{1}(E)$ 的任意 $l$ 个元 $\varphi_{1}\cdots\varphi_{l}$ 生成的 $A(\mathbf{R}^{3})$ 的理想.
\end{definition}

于是,在下文中,关于 $\mathbf{R}^{3}$ 的子流形 $M$ 的具有 restricted cone property(限制锥性质)的紧子集是否为谱综合集,成立如下结果.

\begin{theorem}\label{thm:3.4.25}
设 $M$ 是 $\mathbf{R}^{3}$ 的 $k$ 维子流形($0 \leq k \leq 2$),relative nullity(相对零化度)恒为 $\nu$($0 \leq \nu \leq 2$).又设 $E \subset M$ 是具有 restricted cone property(限制锥性质)的紧集.此时,令 $d(E) = \min\{l \in \mathbb{N}: J_{1}(E)^{l} = J(E)\}$,则 $d(E) \leq \left[k - \frac{\nu}{2} + 1\right]$.特别地,当 $M$ 是超曲面时,$d(E) = \left[k - \frac{\nu}{2} + 1\right]$.
\end{theorem}

关于谱综合,利用定理~\ref{thm:3.4.25} 可证明作为 L. Schwartz 结果的推广的下面的事实.另外,这类讨论在 D. Müller 的论文 \cite{ch3:35} 中对高维情形作了研究.

\begin{theorem}\label{thm:3.4.26}
在定理~\ref{thm:3.4.25} 的条件下,进一步设 $M$ 是 2 维流形.此时,$E$ 是 $A(\mathbf{R}^{3})$ 的谱综合集的充要条件是 $\nu \geq 1$.特别地,$S^{2}$ 不是 $A(\mathbf{R}^{3})$ 的谱综合集.
\end{theorem}
